\section{Giai đoạn Triển khai}

Chương này mô tả chiến lược triển khai và thiết lập hạ tầng cho hệ thống. Hệ thống được phân phối trên các nhà cung cấp dịch vụ đám mây khác nhau nhằm tận dụng tối đa thế mạnh riêng của từng nền tảng: Vercel được sử dụng cho giao diện frontend và backend Express nhờ khả năng tích hợp CI/CD mượt mà và kiến trúc serverless, trong khi Google Cloud Run được sử dụng cho dịch vụ AI để xử lý hiệu quả các tác vụ máy học được đóng gói dạng container.

\subsection{Triển khai Frontend}
Ứng dụng frontend, được xây dựng bằng React và Vite, được triển khai trên Vercel. Vercel cung cấp Mạng lưới Edge Toàn cầu (Edge Network) đảm bảo phân phối nội dung nhanh chóng và tính khả dụng cao. Quy trình triển khai được tích hợp trực tiếp với kho lưu trữ GitHub. Mỗi khi có thao tác push lên nhánh chính (main branch), Vercel sẽ tự động kích hoạt quy trình build mới, thực thi kịch bản biên dịch sản phẩm (production build) và cập nhật trang web trực tuyến. Các biến môi trường, chẳng hạn như URL API backend, được quản lý an toàn trong bảng điều khiển Vercel.

\subsection{Triển khai Backend}
Dịch vụ backend, được xây dựng với Node.js và Express, cũng được triển khai trên Vercel bằng cách sử dụng các Hàm Serverless (Serverless Functions). Tiếp cận này loại bỏ nhu cầu tự quản trị máy chủ thủ công và tự động mở rộng quy mô dựa trên lưu lượng truy cập thực tế.
Để cấu hình triển khai, tệp \texttt{vercel.json} được sử dụng tại thư mục gốc của backend để chỉ định môi trường build và các quy tắc định tuyến. Cấu hình điều hướng tất cả các yêu cầu API đến điểm đầu vào chính \texttt{src/app.ts}, sử dụng môi trường thực thi \texttt{@vercel/node}:

\begin{verbatim}
{
  "version": 2,
  "builds": [
    {
      "src": "src/app.ts",
      "use": "@vercel/node"
    }
  ],
  "routes": [
    {
      "src": "/(.*)",
      "dest": "src/app.ts"
    }
  ]
}
\end{verbatim}
Cấu hình này cho phép việc định tuyến của Express hoạt động mượt mà trong môi trường serverless.

\subsection{Triển khai Dịch vụ AI}
Dịch vụ AI, chịu trách nhiệm xử lý các tác vụ tiêu tốn tài nguyên như xử lý ngôn ngữ tự nhiên và tìm kiếm ngữ nghĩa, được phát triển bằng Python. Dịch vụ này yêu cầu môi trường thực thi riêng biệt với các phụ thuộc phức tạp được định nghĩa trong \texttt{requirements.txt}. Để đảm bảo tính nhất quán giữa các môi trường, dịch vụ được đóng gói bằng Docker và triển khai lên Google Cloud Run.

\begin{itemize}
    \item \textbf{Đóng gói Container:} Tệp \texttt{Dockerfile} được sử dụng để đóng gói ứng dụng Python, các phụ thuộc hệ thống và các gói thư viện Python thành một hình ảnh Docker (Docker image) tiêu chuẩn.
    \item \textbf{Google Cloud Run:} Cloud Run là một nền tảng tính toán được quản lý hoàn toàn, tự động mở rộng quy mô các container không lưu trạng thái. Nền tảng này được lựa chọn cho dịch vụ AI vì có thể cung cấp hiệu quả các tài nguyên CPU và bộ nhớ cần thiết cho các yêu cầu xử lý và tự động thu hẹp về 0 khi không có truy cập, giúp tối ưu hóa chi phí vận hành.
    \item \textbf{Quy trình Triển khai:} Hình ảnh Docker được biên dịch và đẩy lên Google Artifact Registry. Một dịch vụ Cloud Run sau đó được triển khai tham chiếu đến hình ảnh đã tải lên. Dịch vụ được cấu hình an toàn với các biến môi trường cần thiết và được cấp phát đủ tài nguyên để xử lý các tải công việc AI đồng thời.
\end{itemize}

\subsection{Tích hợp Hệ thống}
Kiến trúc triển khai độc lập đảm bảo rằng mỗi dịch vụ có thể mở rộng quy mô một cách riêng biệt dựa trên tải thực tế. Frontend giao tiếp với backend Express trên Vercel thông qua các RESTful API bảo mật chuẩn HTTPS. Backend Express, đến lượt mình, sẽ chuyển tiếp các yêu cầu AI chuyên biệt đến điểm cuối Google Cloud Run, đóng vai trò như một API gateway. Chiến lược triển khai mô-đun này giúp tăng cường khả năng chịu lỗi, tính dễ bảo trì và khả năng mở rộng của toàn bộ hệ thống.